\section*{第 \ref{ch:g10:coordgeom} 章——坐标几何}

\begin{solution}{exo:g10:coordgeom:1}
\omterm{prop:g10:coordgeom:midpoint}{中点}：$\left(\frac{2 + (-4)}{2}, \frac{5 + 1}{2}\right) = (-1, 3)$。
距离：
\[
AB = \sqrt{(-4 - 2)^2 + (1 - 5)^2} = \sqrt{36 + 16} = \sqrt{52}
= 2\sqrt{13}.
\]
\end{solution}

\begin{solution}{exo:g10:coordgeom:2}
$I$ 为 $[AB]$ 的\omterm{prop:g10:coordgeom:midpoint}{中点}，即 $\frac{3 + x_B}{2} = 1$ 且
$\frac{-2 + y_B}{2} = 2$，于是 $x_B = -1$，$y_B = 6$，即 $B(-1, 6)$。
\end{solution}

\begin{solution}{exo:g10:coordgeom:3}
$OP = \sqrt{16 + 1} = \sqrt{17}$; $OQ = \sqrt{9 + 4} = \sqrt{13}$;
$OR = \sqrt{0 + 25} = 5 = \sqrt{25}$。因 $13 < 17 < 25$，
且平方根\omterm{def:g10:functions:function}{函数}\omterm{def:g10:functions:variations}{递增}，故 $Q$ 离原点最近。
\end{solution}

\begin{solution}{exo:g10:coordgeom:4}
边长的平方：
\begin{align*}
AB^2 &= (5-1)^2 + (4-2)^2 = 16 + 4 = 20, \\
BC^2 &= (7-5)^2 + (0-4)^2 = 4 + 16 = 20, \\
AC^2 &= (7-1)^2 + (0-2)^2 = 36 + 4 = 40 .
\end{align*}
$AB = BC$，故三角形在 $B$ 处为等腰三角形。此外
$AB^2 + BC^2 = 40 = AC^2$，由勾股定理的逆定理，
三角形在 $B$ 处也是直角三角形。
\end{solution}

\begin{solution}{exo:g10:coordgeom:5}
\emph{1.} $[AC]$ 的\omterm{prop:g10:coordgeom:midpoint}{中点}为
$\left(\frac{-2+4}{2}, \frac{1+1}{2}\right) = (1, 1)$；$[BD]$ 的\omterm{prop:g10:coordgeom:midpoint}{中点}为
$\left(\frac{1+1}{2}, \frac{3+(-1)}{2}\right) = (1, 1)$。两\omterm{prop:g10:coordgeom:midpoint}{中点}相同，
故 $ABCD$ 为平行四边形。

\emph{2.} $AB^2 = 3^2 + 2^2 = 13$，$BC^2 = 3^2 + (-4)^2 = 25$，
则 $AB \neq BC$，不是菱形。

\emph{3.} $AC^2 = 6^2 + 0^2 = 36$，$BD^2 = 0^2 + (-4)^2 = 16$，
则 $AC \neq BD$，不是矩形。$ABCD$ 只是一般的平行四边形。
\end{solution}

\begin{solution}{exo:g10:coordgeom:6}
计算各点到 $\Omega(2,1)$ 的距离平方，并与 $5^2 = 25$ 比较：
\[
\Omega A^2 = 3^2 + 4^2 = 25, \qquad
\Omega B^2 = (-3)^2 + (-4)^2 = 25, \qquad
\Omega C^2 = 4^2 + 1^2 = 17 .
\]
$A$、$B$ 在圆上；$C$ 在圆内（$17 < 25$）。
\end{solution}

\begin{solution}{exo:g10:coordgeom:7}
$MA^2 = (x - 1)^2 + 1$，$MB^2 = (x - 7)^2 + 25$。令
$MA^2 = MB^2$：
\[
x^2 - 2x + 2 = x^2 - 14x + 74
\ \Longleftrightarrow\
12x = 72
\ \Longleftrightarrow\
x = 6 .
\]
满足条件的点唯一，为 $M(6, 0)$。
\end{solution}

\begin{solution}{exo:g10:coordgeom:8}
\emph{1.} $I = \left(\frac{0+4}{2}, \frac{3+1}{2}\right) = (2, 2)$.

\emph{2.} $M(2,2)$ 确为 $I$。
又 $MA^2 = (0-2)^2 + (3-2)^2 = 5$，$MB^2 = (4-2)^2 + (1-2)^2 = 5$，
可见\omterm{prop:g10:coordgeom:midpoint}{中点}到两端点距离相等，与预期相符。

\emph{3.} 求 $N(0, y)$，使 $NA^2 = NB^2$：
\[
(3 - y)^2 = 16 + (1 - y)^2
\ \Longleftrightarrow\
9 - 6y = 16 + 1 - 2y
\ \Longleftrightarrow\
-4y = 8,
\]
于是 $y = -2$，即点 $N(0, -2)$。
\end{solution}

\begin{solution}{exo:g10:coordgeom:9}
\begin{align*}
AB &= \sqrt{4^2 + 2^2} = \sqrt{20} = 2\sqrt5, \\
BC &= \sqrt{8^2 + 4^2} = \sqrt{80} = 4\sqrt5, \\
AC &= \sqrt{12^2 + 6^2} = \sqrt{180} = 6\sqrt5 .
\end{align*}
于是 $AB + BC = 2\sqrt5 + 4\sqrt5 = 6\sqrt5 = AC$，
故三角形不等式取等号，$B$ 位于线段 $[AC]$ 上——
三点共线。
\end{solution}

\begin{solution}{exo:g10:coordgeom:10}
$[AB]$ 的\omterm{prop:g10:coordgeom:midpoint}{中点}为 $I\left(\frac a2, \frac b2\right)$。于是
\[
OI^2 = \frac{a^2}{4} + \frac{b^2}{4},
\qquad
IA^2 = \left(a - \frac a2\right)^2 + \left(0 - \frac b2\right)^2
= \frac{a^2}{4} + \frac{b^2}{4},
\]
同理 $IB^2 = \left(0 - \frac a2\right)^2 + \left(b - \frac b2\right)^2 =
\frac{a^2}{4} + \frac{b^2}{4}$。三者的平方距离相等，
故 $OI = IA = IB$，即直角三角形 $OAB$ 斜边的\omterm{prop:g10:coordgeom:midpoint}{中点}
与三个顶点等距（它是
其外接圆的圆心）。
\end{solution}

\begin{solution}{pb:g10:coordgeom:1}
\textbf{1.} $\Omega M = 5 \iff \Omega M^2 = 25 \iff
(x - 2)^2 + (y - 1)^2 = 25$ 由距离公式
（\cref{thm:g10:coordgeom:distance}）——平方无妨，
因两边均为正。

\textbf{2.} $A$：$(5-2)^2 + (5-1)^2 = 9 + 16 = 25$，在
圆上。$B$：$16 + 9 = 25$，同样在圆上。$D$：$4 + 9 = 13 < 25$：
在圆内。

\textbf{3.} $x^2 - 6x = (x - 3)^2 - 9$ 及
$y^2 + 4y = (y + 2)^2 - 4$，于是\omterm{def:g10:algebra:equation}{方程}化为
$(x - 3)^2 + (y + 2)^2 = 25$：圆心 $(3, -2)$，半径 $5$。

\textbf{4.} $(x - 3)^2 + (y + 2)^2 = -1$：平方和
永不为负——没有点满足它：空集。一般地
$(x - h)^2 + (y - k)^2 = c$ 是半径为
$\sqrt c$ 的圆（当 $c > 0$ 时），单点 $(h, k)$（当 $c = 0$ 时），而
为空集（当 $c < 0$ 时）。

\textbf{5.} 取 $y = 3$，$(x - 3)^2 + 25 = 25$，于是
$(x - 3)^2 = 0$，得单点 $(3, 3)$。只有一个切点：
直线是圆的\emph{切线}。

\textbf{6.} $MA^2 + MB^2 = (x + r)^2 + y^2 + (x - r)^2 + y^2
= 2x^2 + 2y^2 + 2r^2 = 2r^2 + 2r^2 = 4r^2$（利用
$x^2 + y^2 = r^2$）。由于 $AB^2 = (2r)^2 = 4r^2$，三角形
$AMB$ 满足 $MA^2 + MB^2 = AB^2$，在 $M$ 处为直角
（毕达哥拉斯定理的逆定理）——三行即可，正如所许诺。

\textbf{7.} $MA^2 + MB^2 = (x + a)^2 + y^2 + (x - a)^2 + y^2
= 2x^2 + 2y^2 + 2a^2 = 2\,MI^2 + 2a^2$，又
$\frac{AB^2}{2} = \frac{(2a)^2}{2} = 2a^2$。恒等式得证。

\textbf{8.} 由 $2\,MI^2 + \frac{36}{2} = 26$ 得 $MI^2 = 4$，即
以\omterm{prop:g10:coordgeom:midpoint}{中点} $I$ 为圆心、$2$ 为半径的圆。

\textbf{9.} $MA^2 = 4\,MB^2$：
$x^2 + y^2 = 4\left((x - 3)^2 + y^2\right)$，即\
$3x^2 + 3y^2 - 24x + 36 = 0$，或
$x^2 + y^2 - 8x + 12 = 0$：配方得
$(x - 4)^2 + y^2 = 4$——圆心为 $(4, 0)$，
半径为 $2$ 的\emph{阿波罗尼斯圆}（比值为 $2$）。

\textbf{10.} 这座桥梁把作图转化为
计算：无需辅助矩形，无需分类讨论——一次\omterm{def:g10:algebra:expand}{展开}
便可自然得出定理。代价是：计算
较少体现几何直观；经典证明说明
\emph{为何}，代数证明只说明\emph{如此}——一位
从事研究的数学家两者都会保留。

\textbf{11.} $x^2 + y^2 = 25$ 与 $(x - 6)^2 + y^2 = 25$。
两式相减：$x^2 - (x - 6)^2 = 0$，即\ $12x - 36 = 0$：
$x = 3$。于是 $9 + y^2 = 25$，$y = \pm 4$。候选点为
$(3, 4)$ 与 $(3, -4)$。

\textbf{12.} 自 $(3, 4)$ 至 $R(0, 8)$：
$\sqrt{9 + 16} = 5$——相符。自 $(3, -4)$ 出发：
$\sqrt{9 + 144} \approx 12.4$——不符，故你位于 $(3, 4)$。

\textbf{13.} 两\omterm{def:g10:algebra:equation}{方程}都含相同的 $x^2 + y^2$（或\omterm{def:g10:algebra:expand}{展开}后系数相同），
相减消去平方项，
得一次\omterm{def:g10:algebra:equation}{方程}——即一条直线。当
两圆相交于两点时，该直线经过两个交点，
即公共弦所在的直线。

\textbf{14.} 未知数有四个：三个\omterm{def:g10:coordgeom:system}{坐标}加上
接收机的钟差（一块廉价手表，经
数学修正）。每颗卫星提供一个\omterm{def:g10:algebra:equation}{方程}，因此四颗卫星
给出四个\omterm{def:g10:algebra:equation}{方程}，解四个未知数——多于四颗卫星
时可提高精度。

\textbf{15.} 现相减得
$x = \frac{5.1^2 - 5^2 + 36}{12} = \frac{37.01}{12} \approx
3.084$：某一距离上 $10$ 厘米的误差，使估计值
在此约移动 $8$ 厘米。误差会沿每个公式传播——
即\cref{pb:g10:numbers:1}中\omterm{def:g10:numbers:interval}{区间}运算的警示，
现在则与\omterm{def:g10:coordgeom:system}{坐标}相结合。

\textbf{16.} $B' (7, -3)$。直线 $(AB')$ 的\omterm{def:g10:reffunc:affine}{斜率}为
$\frac{-3 - 5}{7 - 1} = -\frac43$，与 $y = 0$ 相交于
$x = 1 + \frac{5}{8} \times 6 = 4.75$ 处，饮水点为
$P(4.75,\ 0)$。最短路程：$AP + PB = AP + PB' = AB' =
\sqrt{6^2 + 8^2} = 10$。

\textbf{17.} 经由 $Q(3, 0)$：
$AQ + QB = \sqrt{4 + 25} + \sqrt{16 + 9} = \sqrt{29} + 5
\approx 10.39 > 10$。经反射所得的直线路径更短。

\textbf{18.} 对角线\omterm{prop:g10:coordgeom:midpoint}{中点}：$[AC]$：$(3, 3)$；$[BD]$：
$(3, 3)$——相等，故为平行四边形。边长：
$AB = \sqrt{10} = AD$，为菱形。对角线：
$AC = \sqrt{32} \neq BD = \sqrt{8}$，不是矩形，因此
也不是正方形。结论是：仅为菱形，别无其他。

\textbf{19.} $D = A + C - B = (-1 + 6 - 3,\ 2 + 0 - 4) =
(2, -2)$：于是 $[AC]$ 与 $[BD]$ 的\omterm{prop:g10:coordgeom:midpoint}{中点}为
$\left(\frac52, 1\right)$。

\textbf{20.} \omterm{prop:g10:coordgeom:midpoint}{中点}公式：用于平行四边形与中心
（第 18、19 题）。距离公式：圆、轨迹、三角形
的形状、最短路径（第 1--9、16 题）。圆\omterm{def:g10:algebra:equation}{方程}相减：
确定位置的直线（第
11--15 题）。工具箱中缺少的是一套清晰的代数，用于处理
\emph{方向}——即向量——与\emph{\omterm{def:g10:reffunc:affine}{斜率}}——
即直线\omterm{def:g10:algebra:equation}{方程}与\omterm{def:g10:algebra:equation}{方程}组。接下来的两章
恰好补上这一部分。
\end{solution}
